\documentclass[a4paper,12pt]{article}
\usepackage[ngerman]{babel}
\usepackage{amsfonts}
\usepackage{amsmath}
\usepackage{tikz}
\usetikzlibrary{trees}

\title{Übungsblatt 6 - Algorithmen und Datenstrukturen}
\author{Alexander Engelhardt, Timothy Drexler}
\date{\today}

\begin{document}
\maketitle

\section*{Aufgabe 9 (AVL-Bäume)}

\subsection*{Ausgangssituation}
\begin{center}
\begin{tikzpicture}[
  level distance=1.5cm,
  level 1/.style={sibling distance=4cm},
  level 2/.style={sibling distance=2cm},
  level 3/.style={sibling distance=1cm},
  every node/.style={circle,draw,minimum size=0.8cm}
]
\node {10}
  child {node {5}
    child {node {3}}
    child {node {7}
      child {node {6}}
      child[missing]
    }
  }
  child {node {12}
    child {node {11}}
    child {node {15}}
  };
\end{tikzpicture}
\end{center}

\subsection*{Schritt 1: Einfügen von 9}

9 wird als rechtes Kind von 7 eingefügt. Dies verursacht ein Ungleichgewicht bei Knoten 5 (Balance-Faktor = -2).
Es wird eine Links-Rechts-Rotation an Knoten 5 durchgeführt.

\begin{center}
\begin{tikzpicture}[
  level distance=1.5cm,
  level 1/.style={sibling distance=4cm},
  level 2/.style={sibling distance=2cm},
  level 3/.style={sibling distance=1cm},
  every node/.style={circle,draw,minimum size=0.8cm}
]
\node {10}
  child {node {7}
    child {node {5}
      child {node {3}}
      child {node {6}}
    }
    child {node {9}}
  }
  child {node {12}
    child {node {11}}
    child {node {15}}
  };
\end{tikzpicture}
\end{center}

\textbf{Operation:} Links-Rechts-Rotation an Knoten 5 (zuerst Linksrotation an 7, dann Rechtsrotation an 5).

\subsection*{Schritt 2: Einfügen von 8}

8 wird als linkes Kind von 9 eingefügt. Der Baum bleibt balanciert.

\begin{center}
\begin{tikzpicture}[
  level distance=1.5cm,
  level 1/.style={sibling distance=4cm},
  level 2/.style={sibling distance=2cm},
  level 3/.style={sibling distance=1cm},
  every node/.style={circle,draw,minimum size=0.8cm}
]
\node {10}
  child {node {7}
    child {node {5}
      child {node {3}}
      child {node {6}}
    }
    child {node {9}
      child {node {8}}
      child[missing]
    }
  }
  child {node {12}
    child {node {11}}
    child {node {15}}
  };
\end{tikzpicture}
\end{center}

\textbf{Operation:} Keine Rotation erforderlich, Baum bleibt balanciert.

\subsection*{Schritt 3: Einfügen von 13}

13 wird als linkes Kind von 15 eingefügt. Der Baum bleibt balanciert.

\begin{center}
\begin{tikzpicture}[
  level distance=1.5cm,
  level 1/.style={sibling distance=4cm},
  level 2/.style={sibling distance=2cm},
  level 3/.style={sibling distance=1cm},
  every node/.style={circle,draw,minimum size=0.8cm}
]
\node {10}
  child {node {7}
    child {node {5}
      child {node {3}}
      child {node {6}}
    }
    child {node {9}
      child {node {8}}
      child[missing]
    }
  }
  child {node {12}
    child {node {11}}
    child {node {15}
      child {node {13}}
      child[missing]
    }
  };
\end{tikzpicture}
\end{center}

\textbf{Operation:} Keine Rotation erforderlich, Baum bleibt balanciert.

\subsection*{Schritt 4: Löschen von 10}

Knoten 10 (Wurzel) wird entfernt. Der symmetrische Nachfolger ist 11.
11 ersetzt 10, und der Baum wird rebalanciert.

\begin{center}
\begin{tikzpicture}[
  level distance=1.5cm,
  level 1/.style={sibling distance=4cm},
  level 2/.style={sibling distance=2cm},
  level 3/.style={sibling distance=1cm},
  every node/.style={circle,draw,minimum size=0.8cm}
]
\node {11}
  child {node {7}
    child {node {5}
      child {node {3}}
      child {node {6}}
    }
    child {node {9}
      child {node {8}}
      child[missing]
    }
  }
  child {node {13}
    child {node {12}}
    child {node {15}
    }
  };
\end{tikzpicture}
\end{center}

\textbf{Operation:} Löschen von 10, Durch Drehen ersetzen durch 13.

\subsection*{Endergebnis}

Der finale AVL-Baum nach allen Operationen enthält die Schlüssel: 3, 5, 6, 7, 8, 9, 11, 12, 13, 15.

\newpage

\section*{Aufgabe 10 (Topologische Sortierung)}

Für den gegebenen gerichteten Graphen wird die topologische Sortierung durchgeführt.
Initial werden Knoten mit gleichem Eingangsgrad nach ihrer Nummer sortiert.

\subsection*{Kantenstruktur des Graphen}

\textbf{Eingangsgrade:}
\begin{itemize}
    \item Knoten 0: Eingangsgrad = 0
    \item Knoten 1: Eingangsgrad = 1 (von 0)
    \item Knoten 2: Eingangsgrad = 1 (von 0)
    \item Knoten 3: Eingangsgrad = 1 (von 1)
    \item Knoten 4: Eingangsgrad = 1 (von 3)
    \item Knoten 5: Eingangsgrad = 2 (von 3, 4)
    \item Knoten 6: Eingangsgrad = 1 (von 4)
    \item Knoten 7: Eingangsgrad = 2 (von 5, 6)
    \item Knoten 8: Eingangsgrad = 2 (von 12, 7)
    \item Knoten 9: Eingangsgrad = 1 (von 6)
    \item Knoten 10: Eingangsgrad = 2 (von 7, 8)
    \item Knoten 11: Eingangsgrad = 2 (von 9, 10)
    \item Knoten 12: Eingangsgrad = 1 (von 2)
\end{itemize}

\subsection*{Durchführung der topologischen Sortierung}

{\small
\begin{tabular}{|c|l|l|l|l|}
\hline
\textbf{Schritt} & \textbf{Kandidaten} & \textbf{ts-Feld} & \textbf{InDegree-Feld} & \textbf{Entfernt} \\
\hline
0 & [0] & [-,-,-,-,-,-,-,-,-,-,-,-,-] & [0,1,1,2,1,3,2,0,2,2,2,2,0] & - \\
\hline
1 & [7] & [0,-,-,-,-,-,-,-,-,-,-,-,-] & [0,0,1,2,1,3,2,0,2,2,2,2,0] & 0 \\
\hline
2 & [12] & [0,7,-,-,-,-,-,-,-,-,-,-,-] & [0,0,1,2,1,2,2,0,1,1,1,2,0] & 7 \\
\hline
3 & [1] & [0,7,12,-,-,-,-,-,-,-,-,-,-] & [0,0,0,1,1,2,2,0,1,1,1,2,0] & 12 \\
\hline
4 & [2] & [0,7,12,1,-,-,-,-,-,-,-,-,-] & [0,0,0,0,1,2,2,0,1,1,1,2,0] & 1 \\
\hline
5 & [3] & [0,7,12,1,2,-,-,-,-,-,-,-,-] & [0,0,0,0,0,1,2,0,1,1,1,2,0] & 2 \\
\hline
6 & [4] & [0,7,12,1,2,3,-,-,-,-,-,-,-] & [0,0,0,0,0,1,1,0,1,1,1,2,0] & 3 \\
\hline
7 & [5] & [0,7,12,1,2,3,4,-,-,-,-,-,-] & [0,0,0,0,0,0,0,0,0,1,1,2,0] & 4 \\
\hline
8 & [6] & [0,7,12,1,2,3,4,5,-,-,-,-,-] & [0,0,0,0,0,0,0,0,0,0,1,2,0] & 5 \\
\hline
9 & [8] & [0,7,12,1,2,3,4,5,6,-,-,-,-] & [0,0,0,0,0,0,0,0,0,0,0,2,0] & 6 \\
\hline
10 & [9] & [0,7,12,1,2,3,4,5,6,8,-,-,-] & [0,0,0,0,0,0,0,0,0,0,0,1,0] & 8 \\
\hline
11 & [10] & [0,7,12,1,2,3,4,5,6,8,9,-,-] & [0,0,0,0,0,0,0,0,0,0,0,0,0] & 9 \\
\hline
12 & [11] & [0,7,12,1,2,3,4,5,6,8,9,10,-] & [0,0,0,0,0,0,0,0,0,0,0,0,0] & 10 \\
\hline
13 & [] & [0,7,12,1,2,3,4,5,6,8,9,10,11] & [0,0,0,0,0,0,0,0,0,0,0,0,0] & 11 \\
\hline
\end{tabular}
}
\vspace{1cm}
\end{document}
